\documentclass[tikz,border=4pt]{standalone}
\usepackage{tikz}
\usepackage{polymer-paper-preamble}
% Reusable Panel F apparatus motif. The caller owns labels, force arrows,
% panel framing, and whole-panel composition.
\newcommand{\PanelFFloatingCantilever}[1]{%
  \begin{scope}[shift={(#1)}]
    % Stable anchors for semantic overlays and downstream inspection.
    \coordinate (panelFFixedBoundary) at (0,0.34);
    \coordinate (panelFFloatingCantilever) at (-0.26,0);
    \coordinate (panelFDrivenElectrode) at (1.34,-1.40);
    \coordinate (panelFSourceReturn) at (1.74,0.39);
    \coordinate (panelFTrappedCharge) at (-0.22,-0.56);

    % Fixed mechanical boundary and shallow root jaw.
    \draw[cGray!62!black, line width=0.72pt, line cap=round, line join=round]
      (-0.26,0.34) -- (0.26,0.34);
    \foreach \hx in {-0.22,-0.08,0.06,0.20} {
      \draw[cGray!48!black, line width=0.30pt]
        (\hx,0.34) -- ({\hx+0.07},0.42);
    }
    \draw[cGray!62!black, line width=0.72pt, line cap=round]
      (0,0.34) -- (0,0.14);
    \fill[cGray!12, rounded corners=0.22mm]
      (-0.16,0) rectangle (0.16,0.14);
    \draw[cGray!62!black, line width=0.72pt, rounded corners=0.22mm]
      (-0.16,0) rectangle (0.16,0.14);
    \draw[cGray!56!black, line width=0.34pt] (-0.12,0) -- (-0.12,-0.06);
    \draw[cGray!56!black, line width=0.34pt] (0.12,0) -- (0.12,-0.06);

    % Electrically floating charge-trapping cantilever.
    \shade[top color=cAmber!18, bottom color=cAmber!46, rounded corners=0.5mm]
      (-0.14,0) .. controls (-0.02,-0.74) and (-0.49,-1.54) ..
      (-1.37,-2.04) -- (-1.22,-2.20) .. controls (-0.28,-1.72) and
      (0.24,-0.78) .. (0.14,0) -- cycle;
    \draw[cAmber!82!black, line width=0.62pt, rounded corners=0.5mm]
      (-0.14,0) .. controls (-0.02,-0.74) and (-0.49,-1.54) ..
      (-1.37,-2.04) -- (-1.22,-2.20) .. controls (-0.28,-1.72) and
      (0.24,-0.78) .. (0.14,0) -- cycle;
    \draw[cAmber!60!black, line width=0.28pt, opacity=0.42]
      (-0.06,-0.16) .. controls (-0.02,-0.76) and (-0.56,-1.54) .. (-1.29,-2.06);
    \draw[cAmber!60!black, line width=0.28pt, opacity=0.35]
      (0.04,-0.12) .. controls (0.14,-0.82) and (-0.39,-1.72) .. (-1.16,-2.16);

    % Driven electrode and gap guides.
    \shade[left color=cGray!30, right color=cGray!16]
      (1.34,-2.38) rectangle (1.58,-0.02);
    \draw[cGray!82!black, line width=0.62pt]
      (1.34,-2.38) rectangle (1.58,-0.02);
    \draw[cGray!45!black, line width=0.25pt] (1.39,-2.34) -- (1.39,-0.06);
    \foreach \hy in {-2.22,-1.98,-1.74,-1.50,-1.26,-1.02,-0.78,-0.54,-0.30} {
      \draw[cGray!48!black, line width=0.20pt]
        (1.58,\hy) -- (1.34,{\hy-0.05});
    }
    \foreach \yy/\xstart in {-1.88/-0.96,-1.50/-0.67,-1.12/-0.44,-0.74/-0.22} {
      \draw[cGray!23, line width=0.30pt, dash pattern=on 1.3pt off 1.2pt]
        (\xstart,\yy) -- (1.34,\yy);
    }
    \draw[cGray!19, line width=0.22pt, dash pattern=on 1.2pt off 1.5pt]
      (1.08,-0.42) -- (1.34,-0.42);

    % Driven source lead and grounded source return. No lead touches sample.
    \fill[cBlue!5, rounded corners=0.55mm] (0.76,0.40) rectangle (1.24,0.78);
    \draw[cBlue!64!black, line width=0.38pt, rounded corners=0.55mm]
      (0.76,0.40) rectangle (1.24,0.78);
    \node[font=\sffamily\bfseries\fontsize{4.4}{5.2}\selectfont,
          text=cBlue!70!black] at (1.00,0.59) {HV};
    \fill[cBlue!70!black] (1.24,0.50) circle (0.025);
    \draw[cBlue!70!black, line width=0.36pt, rounded corners=0.3mm]
      (1.24,0.50) -- (1.46,0.50) -- (1.46,-0.02);
    \fill[cBlue!70!black] (1.24,0.68) circle (0.025);
    \draw[cBlue!58!black, line width=0.34pt, rounded corners=0.3mm]
      (1.24,0.68) -- (1.74,0.68) -- (1.74,0.39);
    \draw[cBlue!58!black, line width=0.34pt] (1.61,0.39) -- (1.87,0.39);
    \draw[cBlue!58!black, line width=0.34pt] (1.65,0.33) -- (1.83,0.33);
    \draw[cBlue!58!black, line width=0.34pt] (1.69,0.27) -- (1.79,0.27);

    % Flat trapped-charge markers.
    \foreach \cx/\cy/\rr in {-0.22/-0.56/0.075,-0.41/-0.98/0.082,-0.74/-1.48/0.088,-1.12/-1.90/0.082} {
      \fill[cRed!68!black] (\cx,\cy) circle (\rr);
      \draw[cRed!88!black, line width=0.25pt] (\cx,\cy) circle (\rr);
    }
  \end{scope}%
}


\tikzset{
  panel letter/.style={
    font=\sffamily\bfseries\fontsize{9}{10.8}\selectfont,
    text=black
  },
  panel title/.style={
    font=\sffamily\bfseries\fontsize{7.2}{8.6}\selectfont,
    text=cGray!92!black,
    align=left
  },
  body label/.style={
    font=\sffamily\fontsize{6.0}{7.2}\selectfont,
    text=cGray,
    align=center
  },
  small label/.style={
    font=\sffamily\fontsize{5.2}{6.3}\selectfont,
    text=cGray,
    align=center
  },
  evidence axis/.style={
    -{Stealth[length=3pt,width=2.2pt]},
    draw=cGray,
    line width=0.42pt
  },
  evidence tick/.style={draw=cGray, line width=0.25pt},
  guide/.style={draw=cLGray, line width=0.32pt, dashed},
  relation/.style={
    -{Stealth[length=3.4pt,width=2.5pt]},
    draw=cGray!72,
    line width=0.45pt
  },
  force arrow/.style={
    -{Stealth[length=4.2pt,width=3.2pt]},
    draw=cRed!88!black,
    line width=0.85pt
  }
}

\begin{document}
\begin{tikzpicture}[x=1cm,y=1cm,line cap=round,line join=round]
  % Figure heading and the material-to-evidence narrative spine.
  \node[
    anchor=west,
    font=\sffamily\bfseries\fontsize{11}{13.2}\selectfont,
    text=black
  ] at (0,12.15) {Sulfur-rich polymers localize charge across multiple trap depths};
  \draw[cAmber, line width=1.05pt] (0,11.78) -- (19.0,11.78);

  % A — material identity.
  \node[panel letter, anchor=west] at (0,11.30) {A};
  \node[panel title, anchor=west] at (0.42,11.30) {Sulfur-rich polymer};
  \WavyChain{0.20,4.55,10.20,zigzag,{0.85,1.75,2.70,3.65,4.25},{}}
  \fill[cAmberSphere!34, rounded corners=2pt]
    (0.25,8.92) rectangle (4.52,9.58);
  \draw[cAmber!82!black, line width=0.55pt, rounded corners=2pt]
    (0.25,8.92) rectangle (4.52,9.58);
  \foreach \sx/\sy in {0.62/9.25,1.18/9.12,1.74/9.38,2.32/9.19,2.91/9.34,3.46/9.10,4.08/9.31} {
    \fill[cAmber] (\sx,\sy) circle (0.065);
  }
  \node[body label] at (2.38,8.68) {sulfur-bearing network};

  % B — composition series, deliberately tucked beneath the identity motif.
  \node[panel letter, anchor=west] at (0,8.04) {B};
  \node[panel title, anchor=west] at (0.42,8.04) {Sulfur-content series};
  \foreach \x/\name in {0.35/low,1.82/intermediate,3.29/high} {
    \fill[cLGray!28, rounded corners=1.5pt]
      (\x,6.62) rectangle ++(1.02,0.78);
    \draw[cGray!76!black, line width=0.45pt, rounded corners=1.5pt]
      (\x,6.62) rectangle ++(1.02,0.78);
    \node[small label] at (\x+0.51,6.38) {\name};
  }
  \foreach \x/\y in {0.72/6.94,2.09/6.86,2.54/7.16,3.49/6.82,3.79/7.13,4.06/6.92,3.65/7.27} {
    \fill[cAmber] (\x,\y) circle (0.075);
  }
  \draw[relation] (0.56,6.02) -- (4.18,6.02);
  \node[small label, fill=white, inner sep=1pt] at (2.37,6.02)
    {increasing sulfur content};

  % C — central qualitative energy landscape.
  \node[panel letter, anchor=west] at (5.22,11.30) {C};
  \node[panel title, anchor=west] at (5.64,11.30)
    {Localized shallow and deep traps coexist};
  \draw[evidence axis] (5.70,6.34) -- (5.70,10.72);
  \node[body label, rotate=90] at (5.35,8.53) {Energy};
  \draw[guide] (6.10,10.16) -- (11.36,10.16);
  \node[small label, anchor=west] at (11.43,10.16) {transport level};
  \draw[cGray!75!black, line width=0.65pt]
    (6.14,9.58)
    .. controls (6.72,9.56) and (6.82,8.78) .. (7.42,8.78)
    .. controls (8.02,8.78) and (8.13,9.55) .. (8.73,9.58)
    .. controls (9.16,9.58) and (9.24,7.33) .. (9.96,7.33)
    .. controls (10.70,7.33) and (10.78,9.57) .. (11.34,9.58);
  \fill[cBlue!65] (7.42,8.95) circle (0.105);
  \draw[cBlue!86!black, line width=0.30pt] (7.42,8.95) circle (0.105);
  \fill[cRed!70!black] (9.96,7.50) circle (0.105);
  \draw[cRed!90!black, line width=0.30pt] (9.96,7.50) circle (0.105);
  \draw[relation] (7.42,8.66) -- ++(0,-0.76);
  \draw[relation] (9.96,7.20) -- ++(0,-0.76);
  \node[body label, text=cBlue!80!black] at (7.42,7.60) {shallow};
  \node[body label, text=cRed!88!black] at (9.96,6.16) {deep};
  \node[small label, text=cGray!85!black] at (8.70,5.72)
    {qualitative localized-state landscape};

  % D — measured transient-current evidence.
  \node[panel letter, anchor=west] at (12.55,11.30) {D};
  \node[panel title, anchor=west] at (12.97,11.30)
    {Transient-current kinetics};
  \draw[evidence axis] (13.05,8.18) -- (18.62,8.18);
  \draw[evidence axis] (13.05,8.18) -- (13.05,10.64);
  \node[body label] at (15.83,7.77) {time};
  \node[body label, rotate=90] at (12.61,9.42) {current};
  \node[small label, anchor=east, text=cBlue!80!black] at (18.58,10.38)
    {constant voltage};
  \draw[cBlue, line width=0.92pt]
    (13.26,10.34)
    .. controls (13.72,9.35) and (14.65,8.88) .. (15.56,8.62)
    .. controls (16.48,8.38) and (17.56,8.30) .. (18.42,8.27);
  \foreach \x/\y in {13.26/10.34,13.72/9.53,14.38/9.06,15.18/8.70,16.18/8.45,17.22/8.33,18.42/8.27} {
    \fill[white] (\x,\y) circle (0.070);
    \draw[cBlue, line width=0.38pt] (\x,\y) circle (0.070);
  }
  \node[small label, anchor=west] at (14.26,10.04)
    {trap-mediated decay};

  % E — noncontact ISPD observation and its derived distribution.
  \node[panel letter, anchor=west] at (12.55,7.27) {E};
  \node[panel title, anchor=west] at (12.97,7.27)
    {ISPD-derived trap distribution};
  \fill[cAmberSphere!30, rounded corners=1.5pt]
    (12.95,5.48) rectangle (15.42,5.92);
  \draw[cAmber!84!black, line width=0.48pt, rounded corners=1.5pt]
    (12.95,5.48) rectangle (15.42,5.92);
  \foreach \x in {13.25,13.68,14.11,14.54,14.97} {
    \fill[cRed!76!black] (\x,5.72) circle (0.055);
  }
  \draw[cGray!78!black, line width=0.55pt, fill=cLGray!24]
    (13.72,6.75) -- (14.66,6.75) -- (14.45,6.24) -- (13.93,6.24) -- cycle;
  \draw[guide] (14.19,6.22) -- (14.19,5.96);
  \node[small label] at (14.18,5.18) {noncontact surface-potential decay};
  \draw[relation] (15.66,5.70) -- (16.26,5.70);
  \draw[evidence axis] (16.52,4.67) -- (18.72,4.67);
  \draw[evidence axis] (16.52,4.67) -- (16.52,6.62);
  \node[small label] at (17.63,4.34) {trap depth};
  \node[small label, rotate=90] at (16.15,5.65) {distribution};
  \path[draw=cTeal, fill=cTeal!18, line width=0.70pt]
    (16.68,4.68)
    .. controls (16.95,4.72) and (16.96,5.48) .. (17.26,5.54)
    .. controls (17.55,5.49) and (17.57,4.78) .. (17.82,4.72)
    .. controls (18.02,4.78) and (18.10,6.08) .. (18.38,6.16)
    .. controls (18.63,6.02) and (18.62,4.77) .. (18.72,4.68) -- cycle;
  \node[small label, text=cTeal!72!black] at (17.25,5.86) {shallow};
  \node[small label, text=cTeal!72!black] at (18.35,6.48) {deep};

  % F — curated floating-cantilever apparatus, composed as the lower anchor.
  \node[panel letter, anchor=west] at (5.22,4.76) {F};
  \node[panel title, anchor=west] at (5.64,4.76)
    {Floating cantilever Coulomb response};
  \PanelFFloatingCantilever{pf}{(8.34,3.82)}
  \draw[force arrow]
    ($(pfTrappedCharge)+(-0.26,-0.48)$) -- ++(-1.28,-0.38);
  \node[body label, anchor=east, text=cRed!88!black]
    at ($(pfTrappedCharge)+(-1.58,-0.90)$) {Coulomb force};
  \draw[relation]
    ($(pfFloatingCantilever)+(-0.18,-0.12)$) -- ++(-1.16,0.20);
  \node[body label, anchor=east]
    at ($(pfFloatingCantilever)+(-1.48,0.10)$) {electrically floating};
  \draw[relation]
    ($(pfSourceReturn)+(0.14,0.10)$) -- ++(0.82,0.36);
  \node[small label, anchor=west]
    at ($(pfSourceReturn)+(1.02,0.49)$) {grounded source return};
  \node[small label, anchor=west, text=cBlue!78!black]
    at ($(pfDrivenElectrode)+(0.46,-0.26)$) {driven electrode};

  % Compact synthesis statement, not a quantitative causal law.
  \draw[cLGray, line width=0.40pt] (0.18,0.62) -- (18.82,0.62);
  \node[
    font=\sffamily\fontsize{6.2}{7.4}\selectfont,
    text=cGray!92!black,
    align=center
  ] at (9.50,0.22)
    {Composition-dependent localized traps are supported by kinetic and noncontact electrical evidence, then expressed mechanically by electrostatic force.};
\end{tikzpicture}
\end{document}
